\documentclass{article}
\usepackage{amsmath,hyperref}
\begin{document}
\section*{A strict arclength separation for early information in the distal cave}
\paragraph{Held-in operational model.}
Retain Cave Circuit 01 placements, the seven pinned collision
meshes, the 20-tile collision scene and the same authorized
continuous center domain: the union of closed one-meter cubes
from the original lookup at these tiles. In the two alternative
worlds exactly one of the two authored portal slabs is closed.
The robot is an axis-aligned cube of half-width 0.3 m.
Start from the checked pose in tile 103, visit the specified
inspection pose in 132, and return to the original pose. The
criterion here is arclength, not acceleration, braking, sensing
latency, packet reliability, or an independent mission. The
point visit is the modeled service; no physical sensing at that
pose or report delivery is credited by visiting it.

At early reveal the closure bit is available before departing.
Under late-only information no status distinguishes the worlds
before entering one of the named gate endpoint tiles; allow
instant correct revelation at that first entry. No assumption
is made that later motion must follow a waypoint library.
That local observation rule and center domain are declared:
earlier physical perception, motion outside the domain, or a
different fault process need another test.

\paragraph{Executable upper bound.}
Both early branches have a saved round-trip polyline with
exact verification. A translating fixed cube sweeps the
Minkowski sum of the segment with the cube, not its entire
axis-aligned bounding box. Candidate separating axes were
found numerically, but each separation is then checked with
integer projections against every triangle surviving an
exact integer bounding-box broad phase; the 0.5 mm yaw pad
is included. The slab is checked with exact fractional segment
intersection against the inflated box. All center chords
remain in the union of allowed cubes. The exact rounded-up
arclengths are 393.522984 and 444.274308 meters. Reproducing
these upper bounds does not require reproducing the path search.
Both saved paths also pass with a 0.302-meter padded half-width;
this check does not imply centimeter-level tolerance.

\paragraph{A lower bound against every late-only motion.}
Use every source cube-face, edge and corner contact. Enclose
each inter-tile contact in its checked portal box. For any
continuous path, an ordered subsequence of such contacts
induces a simple tile route in a relaxed model in which motion
within a tile is unconstrained. For example, take the last
visit to the current tile along the relevant curve prefix;
the next tile that the curve occupies along a sequence of
times decreasing to that cut shares the contact, and has a
strictly later last visit. The finite last-visit sequence
terminates at the goal. Test this separately on the shared
prefix, the remainder after the first revelation, and the
return from the inspection pose. This does not erase a
required observation: the shared prefix is stopped at its
first observation and no later tile is used to identify it.

No unrevealed route can pass a portal closed in one world
or reach the inspection. There are two possible simple
first-reveal prefixes, starting on the two branches; all
simple post-reveal routes are checked in both worlds. Waiting
or oscillating before that cut does not reduce arclength.
For each fixed portal-box sequence $(B_j)$, use a saved
vector $u_j$ of norm at most one for each inter-box segment.
For arbitrary actual crossing points $x_j\in B_j$,
\[
 \sum_j \|x_{j+1}-x_j\|\ \ge\
 \sum_j u_j\mathbin{\cdot}(x_{j+1}-x_j)\ \ge\
 \min_{x_j\in B_j}\sum_j u_j\mathbin{\cdot}(x_{j+1}-x_j).
\]
Compute the last minimum at the proper box endpoints with
exact integers. A numerical optimizer only proposed the
vectors, and no claim of solving the convex problem exactly
is needed. The certificate rechecks all 12 sequences, including
the tied prefix and outbound remainder per world. Combining
the per-world bounds after each possible first reveal gives
a worst-case lower bound of 456.0981695 meters for every
late-only policy, with the opposing world chosen against
its shared first action.

Thus the certified arclength gap is at least 11.8238615 m
between a feasible early-information policy and every
late-only policy in this model. This is a geometric and
informational discriminator, not an advantage against an
equally informed planner. A deadline inside the gap can be used
for a later integrated \emph{development} test. Neither timely
production of the bit nor compound-fault settlement has been
established here. It does not spend the independent cave mission
or establish I2 acceptance.
\end{document}
